%%%%%%%%%%%%%%%%%%%%%%%%%%%%%%%%%%%%%%%%%%%%%%%%%%%%%%%%%%%%%%%%%%%%%%
% Lớp 11 | Lý thuyết đồ thị | Đường đi Euler và Harmilton
% Dạng 1C2?2-2: Đường đi Harmilton
% Số câu: 1  (sinh tự động bởi tach_cau.py; từ nay thầy sửa trực tiếp file này)
%%%%%%%%%%%%%%%%%%%%%%%%%%%%%%%%%%%%%%%%%%%%%%%%%%%%%%%%%%%%%%%%%%%%%%

%%%%%%%%%%%%%%%% TRẮC NGHIỆM 4 PHƯƠNG ÁN %%%%%%%%%%%%%%%%

% nguồn: Ok_TracNghiemLop11_Lan1
\begin{ex}%[1C2Y2-2]%[ma:ca217ece]
	Cho cấp số cộng $\left(u_n\right)$ có số hạng đầu $u_1=-5$, công sai $d=4$. Công thức của số hạng  tổng quát $u_n$ là:
	 \choice
		 {$u_n=-5+4n$}
		 {$u_n=-1-4n$}
		{$u_n=-5+4n^2$} 
		{\True  $u_n=-9+4n$}
	\loigiai{ 
		Với mọi số  tự nhiên $n \ge 2$ ta có $u_n=u_1+(n-1)d=-5+(n-1).4=4n-9$.
	 }
\end{ex}
